\chapter{Conclusion and Future Work}\label{ch:conclusion}

\section{Conclusion}

This project set out to make diffusion models faster to sample without retraining them, by improving training-free feature caching. A cache makes two decisions at every step, \emph{when} to evaluate the network and \emph{what} to write on the steps it skips. The literature review (\cref{ch:literature}) showed that existing work improved these decisions in isolation, forecast the full hidden state rather than the part that is actually unknown, tuned the forecasters' hyper-parameters empirically, and, without exception, left the errors made on skipped steps uncorrected. The project produced two methods that address these gaps.

\paragraph{ReSPect} (\cref{ch:respect}) couples SeaCache's spectrally aligned gate, derived from the linear-MMSE (Wiener) denoiser, with a forecaster of the block-stack \emph{residual}, the only part of a skipped step's features that is not known exactly. The forecaster blends a ridge-regularised Chebyshev fit of the whole residual history with a first-order discrete Taylor step. We showed that copy-reuse has an error bound linear in the skip length, that Taylor extrapolation attains its ideal $(P{+}1)$-order rate only under a node-spacing condition that a dynamic gate does not guarantee (with a sharp counterexample), and that the Chebyshev truncation error is independent of the skip length; we derived the blend weight and the warm-up from bias--variance and degrees-of-freedom arguments. At matched or lower compute, ReSPect achieved speedups of $3.30\times$, $2.63\times$ and $2.77\times$ on FLUX.1-dev, Wan2.1-1.3B and HunyuanVideo, and improved on the strongest compared baseline by up to 3.46, 11.4 and 9.9\,dB PSNR respectively.

\paragraph{KalSea} (\cref{ch:kalsea}) closes the loop that every previous filler leaves open. A damped-trend Kalman filter forecasts the residual on skipped steps, and when the next step is computed, predictor--corrector gap smoothing (PCGS) uses the Rauch--Tung--Striebel smoother of the filter's own model to correct the residuals of the whole gap, re-integrates the gap through the output layers only, and continues sampling from the corrected latent. PCGS adds no transformer evaluations and no hyper-parameters; we derived its gain in closed form and extended it to multistep solvers, classifier-free guidance and two-stream video-and-audio generation. On FLUX.1-dev, KalSea reached 27.72\,dB at 18.1 of 50 evaluations (SeaCache: 26.29\,dB at 18.4) and 22.53\,dB at 13 (SeaCache: 21.33\,dB); PCGS alone contributed 0.81 and 0.61\,dB and also improved reuse and linear extrapolation. On HunyuanVideo it improved on SeaCache in the same setup by 3.17 and 1.92\,dB, and on Wan2.1 it exceeded SeaCache's published results by 5.07 and 2.07\,dB. The released implementation supports five backbones, including FLUX.2 [dev] and MiniMax-H3, and is verified against exact invariants.

\paragraph{Applications in mechanical engineering} (\cref{ch:applications}). We carried both caches over to two open-weight diffusion models from mechanical engineering: structural topology optimisation with Optimize Any Topology and the reconstruction of a flow past a cylinder with DiffusionPDE. This required generalising the residual split to U-Net denoisers, the gate to timestep-modulated convolutional inputs, the SEA coefficients to DDIM and EDM schedules, and PCGS to DDIM with classifier-free guidance and to gradient-guided second-order samplers; all of it was validated against exact invariants before use. On DiffusionPDE, KalSea reduced the network evaluations of the 2000-step guided sampler by $1.7$--$2.6\times$ and was more accurate than ReSPect in all eight forward and inverse comparisons, recovering the flow of the first released case with 3.3\% error (the paper reports $3.9\%$ on average for the uncached sampler). The OAT evaluation follows the original paper's protocol on its full 5000-problem test split, with compliance computed by the model's own finite-element analysis.

Taken together, the project shows that the two decisions of a cache should be designed jointly, that forecasting the residual rather than the hidden state is both theoretically and practically preferable, and that a cache need not be open loop: the observation made at each computed step can be used not only to predict the future but also to repair the past, at almost no cost.

\section{Future work}

\begin{itemize}
  \item \textbf{Extend the engineering evaluation.} Evaluate DiffusionPDE on all 1000 test scenes of each PDE (Darcy, Poisson, Helmholtz, Navier--Stokes) as the scenes become available, and measure clean wall-clock speedups on dedicated GPUs.
  \item \textbf{Guidance-aware gating.} On DiffusionPDE the safe skipping budget depended on the case; a gate that also measures the change of the guidance signal could choose it automatically.
  \item \textbf{Combine the two methods.} PCGS improved every filler on which it was tested; wrapping it around ReSPect's Chebyshev--Taylor forecaster, with gains from a linear interpolation of the innovation, is a natural next experiment.
  \item \textbf{Choose the computed steps jointly with the corrector.} With a closed loop, the gate could account for the fact that the error of a gap will be corrected, for example by using the filter's predicted uncertainty to decide when to compute.
  \item \textbf{Correct caches inside the block stack.} Many caches reuse features of individual blocks, attention or MLP layers; extending smoothing and re-integration to such intermediate features is open.
  \item \textbf{Larger corrections.} PCGS attaches the residual of the uncorrected latent to the corrected one, which is accurate only for small corrections; long gaps could benefit from a cheap first-order update of the residual itself.
  \item \textbf{Guided and physics-constrained samplers.} Engineering samplers often add gradient guidance or constraint projections. Understanding how skipped steps affect such guidance, and how correction interacts with it, would make caching safe to use in design and simulation workflows.
\end{itemize}
